\subsection{Objetivo: estructura empírica sin target}

Checkpoint~03B responde una pregunta distinta de 03A. En lugar de codificar primero una
hipótesis agronómica, construye descriptores matemáticos de la geometría observada de $X$:
curvatura temporal, persistencia, cambio relativo robusto, concordancia de forma entre
sensores y posición de 2025 respecto al historial corto.

El build canónico retiene 335 features en 197 parcelas, sin missing ni infinitos:
91 clean y 244 competition. Las familias son:
\[
162\ \text{temporal shape}
+42\ \text{symmetric change}
+41\ \text{short-baseline condition}
+72\ \text{aggregate geometry}
+18\ \text{sensor shape similarity}.
\]
Sólo se elimina una columna constante. El target no participa en la construcción de esta
tabla.

\subsection{Geometría temporal}

Dada una serie $x_1,\ldots,x_T$, se definen diferencias locales
\[
s_t=\frac{x_t-x_{t-1}}{m_t-m_{t-1}}.
\]
A partir de ellas se construyen:

\paragraph{Variación total}
\begin{equation}
TV(x)=\sum_{t=2}^{T}|x_t-x_{t-1}|.
\label{eq:tv}
\end{equation}

\paragraph{Rugosidad}
\begin{equation}
R(x)=\frac{1}{T-2}\sum_t|s_t-s_{t-1}|.
\label{eq:roughness}
\end{equation}

\paragraph{Energía de curvatura}
\begin{equation}
C(x)=\frac{1}{T-2}\sum_t(s_t-s_{t-1})^2.
\label{eq:curvature}
\end{equation}

La variación total mide movimiento acumulado, mientras que rugosidad y energía de curvatura
penalizan cambios en la pendiente. Dos parcelas pueden tener la misma amplitud y AUC, pero
trayectorias muy distintas en estas métricas.

También se calculan autocorrelación lag-1, $R^2$ de tendencia lineal, fracción de cambios de
signo en las pendientes, sharpness del máximo, centro de masa temporal y entropía de forma.
Para el centro de masa se normaliza primero:
\[
w_t=\frac{x_t-\min x}{\max x-\min x},
\qquad
M(x)=\frac{\sum_t m_t w_t}{\sum_t w_t}.
\]
Para la entropía, si $p_t=w_t/\sum_s w_s$,
\[
H(x)=-\frac{\sum_t p_t\log p_t}{\log T}.
\]
$H$ cercano a uno describe una señal temporal difusa; valores menores concentran masa en
pocos meses.

\subsection{Condición 2025 respecto a un historial corto}

La posición del valor 2025 dentro del rango histórico parcelario se define como
\begin{equation}
Q_{i,m}
=
\frac{x_{i,2025,m}-x_{i,\mathrm{hist,min}}}
{x_{i,\mathrm{hist,max}}-x_{i,\mathrm{hist,min}}+\epsilon}.
\label{eq:hist-position}
\end{equation}

Para NDVI se conserva una versión escalada y nombrada explícitamente
\artifact{vci_like_short}. Está inspirada en la lógica de índices de condición vegetal
\citep{Bokusheva2016,Serban2025}, pero no se presenta como VCI estándar: el historial sólo
abarca 2022--2024 y los extremos son resúmenes parcelarios, no una climatología larga.
Además, los valores no se recortan a $[0,1]$, de modo que una condición 2025 fuera del rango
histórico sigue siendo observable.

\subsection{Cambio simétrico estable}

Los índices espectrales signed pueden acercarse a cero; una razón ordinaria puede explotar.
Se introduce el contraste
\begin{equation}
D(a,b)=\frac{2(a-b)}{|a|+|b|+\epsilon}.
\label{eq:symmetric-change}
\end{equation}
Esta transformación es adimensional, simétrica respecto a la magnitud de ambos términos y
evita el comportamiento inestable de $a/b$ cuando $b\approx0$.

Para una secuencia de cambios mensuales se resumen media estacional, RMS, máximo absoluto,
media tardía, fracción de cambios positivos y fracción de inversiones de signo. Estos
descriptores son empíricos; no se etiquetan como índices agronómicos canónicos.

\subsection{Similitud temporal entre sensores}

Además del error de nivel de 03A, 03B compara \emph{forma} temporal Sentinel-2 versus Planet
mediante Pearson, Spearman, cosine similarity y RMSE normalizado por rango combinado. Sea
$\bm x$ una trayectoria S2 y $\bm z$ una Planet:
\[
\cos(\bm x,\bm z)
=
\frac{\bm x^\top\bm z}
{\|\bm x\|_2\|\bm z\|_2}.
\]
Las métricas no requieren que ambos sensores tengan calibración idéntica para detectar si
suben y bajan de manera concordante.

\subsection{Descubrimiento supervisado separado del dataset materializado}

La contribución metodológica más delicada de 03B es separar la capa $X$-only de un
\emph{miner} supervisado que nunca se materializa globalmente. Se configuran 24 primitivas.
Dentro de cada training fold:
\begin{enumerate}
\item se rankean primitivas usando únicamente $y$ visible en ese training fold;
\item se conservan las 12 primeras;
\item se generan expresiones de una familia predefinida: sumas, diferencias, productos,
safe ratios, reverse ratios y symmetric change;
\item se rankean por asociación Spearman dentro del training fold;
\item se conservan hasta 20;
\item esas expresiones se transforman sobre train/validation sin volver a usar $y$ de
validación.
\end{enumerate}

Formalmente, si $A$ denota el algoritmo de descubrimiento y $f$ un fold,
\[
\mathcal E_f=A(X_{\mathrm{train}(f)},y_{\mathrm{train}(f)}),
\]
y \emph{no}
\[
\mathcal E=A(X_{\Train},y_{\Train})
\]
reutilizado luego en los folds. La segunda construcción filtraría información de los holdouts
hacia la representación.

\subsection{Auditoría de recurrencia}

La auditoría de estabilidad usa 138 parcelas, 24 primitivas, 12 primitivas por fit y 20
expresiones por fit. A través de los diez folds de los dos protocolos se seleccionaron 88
expresiones únicas. El objetivo no era estimar performance, sino identificar si el miner
producía estructuras recurrentes.

Siete expresiones aparecieron al menos en tres de cinco folds bajo \emph{ambos} protocolos.
Sus componentes se concentran en:
\begin{itemize}
\item tendencia y media reciente de rendimiento municipal histórico SIAP;
\item precipitación estacional 2025;
\item GDD proxy base $5^\circ$C;
\item carbono orgánico de suelo 0--30 cm.
\end{itemize}

Ejemplos estructurales:
\[
\frac{P_{2025}}
{|\mathrm{SIAPtrend}|+\epsilon},
\qquad
\frac{SOC_{0-30}}
{|P_{2025}|+\epsilon},
\]
y
\[
\frac{2(\mathrm{SIAPtrend}-P_{2025})}
{|\mathrm{SIAPtrend}|+|P_{2025}|+\epsilon}.
\]

La alta correlación Spearman de entrenamiento de estas fórmulas no es evidencia de
generalización. Sólo 03C, refitando el miner dentro de cada outer fold, puede responder si la
familia de descubrimiento agrega valor predictivo.

\subsection{Resultado conceptual}

03B evita dos errores frecuentes en feature engineering small-$n$:
\begin{enumerate}
\item generar miles de combinaciones globales y escoger las que correlacionan con todo
$y_{\Train}$ antes de CV;
\item confundir estabilidad de una asociación dentro de train con performance out-of-sample.
\end{enumerate}
La capa determinista queda disponible para modelos y embeddings; el descubrimiento supervisado
queda encapsulado como transformer fold-local.
